\chapter{Conclusão e Trabalhos Futuros}
\label{cap:conclusao}

Este capítulo resume as contribuições técnicas e pedagógicas do trabalho, discute suas limitações e aponta caminhos para continuá-lo.

\section{Considerações Finais}
\label{sec:consideracoes-finais}
A plataforma proposta cumpre o objetivo de democratizar o acesso à prática de Sistemas Operacionais: elimina as barreiras da virtualização e oferece ao professor um ambiente robusto, seguro e com métricas precisas de aprendizagem.

\section{Trabalhos Futuros}
\label{sec:trabalhos-futuros}
Como continuidade deste trabalho, propõe-se:
\begin{enumerate}
    \item ampliar o conjunto de comandos POSIX e incluir utilitários de rede, como \texttt{tcpdump}, \texttt{curl} e \texttt{iptables};
    \item integrar a plataforma ao Moodle da UTFPR por meio do protocolo LTI (\textit{Learning Tools Interoperability});
    \item usar técnicas de inteligência artificial para identificar padrões de dificuldade e recomendar desafios de fixação adaptados a cada aluno.
\end{enumerate}
